% Esqueleto do artigo (SBC, 5 a 8 páginas), na estrutura do CLAUDE.md, seção 9.
% [PENDENTE] modelo SBC "trabalhado em sala" (pendência 4). Até lá, use o modelo público da SBC:
% coloque sbc-template.sty e sbc.bst nesta pasta.
% Compilar: pdflatex artigo && bibtex artigo && pdflatex artigo && pdflatex artigo
% Regra 2: nenhum número digitado à mão. Tabelas e figuras vêm de src/report.py.

\documentclass[12pt]{article}
\usepackage{sbc-template}
\usepackage{graphicx,url}
\usepackage[utf8]{inputenc}
\usepackage[brazil]{babel}
\graphicspath{{../results/figures/}}

\sloppy

\title{Análise inteligente de vídeos de aula: transcrição, OCR e indexação multimodal
com avaliação de qualidade e custo computacional}

\author{Autor 1\inst{1}, Autor 2\inst{1}, Autor 3\inst{1}}
\address{Universidade Federal do Pará (UFPA) -- Belém, PA -- Brasil
  \email{\{autor1,autor2,autor3\}@...}
}

\begin{document}
\maketitle

\begin{abstract}
% TODO (inglês)
\end{abstract}

\begin{resumo}
% TODO: problema, proposta, método (WER/CER, RTF, ruído), resultado principal.
\end{resumo}

\section{Introdução}
% Vídeos de aula sem legenda nem indexação. Contribuição. Fronteira análise × transporte
% (plano de conteúdo × plano de entrega). Pergunta de investigação e secundária (F9).

\section{Fundamentação teórica}
% STT: áudio digital (amostragem a 16 kHz), Whisper (Transformer encoder-decoder) \cite{radford2022}, VAD.
% OCR \cite{tesseract}. Embeddings multimodais \cite{radford2021clip}. HLS e QoE.
% Métricas: WER, CER, RTF, precisão e recall de palavras. LLMs e sumarização \cite{lin2004}.

\section{Trabalhos relacionados e tecnologias}
% Whisper, CLIP, legendas automáticas de plataformas; faster-whisper \cite{fasterwhisper}, jiwer \cite{jiwer}.

\section{Metodologia}
% Pipeline, dataset (licenças), gabarito humano, protocolo de ruído (SNR 20/10/0 dB, seed 0),
% protocolo de tempo (aquecimento, 3 repetições, carga à parte), avaliação às cegas da F9.

\section{Arquitetura e implementação}
% Diagrama em dois planos. Por que faster-whisper/CTranslate2 e int8; por que enviar o SRT ao LLM;
% filtro de alucinações (docs/decisoes.md).

\section{Experimentos e métricas}
% Hardware (results/hardware.json). Matriz RQ1, RQ2 e RQ5.

\section{Resultados e discussão}

\begin{figure}[ht]
\centering
\includegraphics[width=.48\textwidth]{fig1_wer_snr.pdf}\hfill
\includegraphics[width=.48\textwidth]{fig2_rtf_wer.pdf}
\caption{Robustez ao ruído (WER × SNR) e trade-off qualidade × custo (RTF × WER).}
\label{fig:stt}
\end{figure}

% Tabela 1: results/tables/tabela1_wer.md. Tabela 3: tabela3_tempo_stt.md.
% Análise de erros (regra 11): nomes próprios, siglas, números, sotaque, ruído.

\begin{figure}[ht]
\centering
\includegraphics[width=.7\textwidth]{fig4_wer_nota_resumo.pdf}
\caption{F9: WER do modelo de STT × nota do resumo gerado por LLM.}
\label{fig:f9}
\end{figure}

\section{Conclusão e trabalhos futuros}
% Limitações: ruído branco sintético, amostra pequena, um orador, não determinismo do LLM.
% Futuro: diarização, ruído real de sala, F6 a F8.

\bibliographystyle{sbc}
\bibliography{referencias}

\end{document}
